% fixture: babel/spanish-report
% tier: T3
% engine: document
% covers: spanish in report with \pagestyle{headings}: \chaptername (Cap\'{\i}tulo), \contentsname (\'{I}ndice general), \bibname (Bibliograf\'{\i}a), the marks in \foreignlanguage{spanish} and their \MakeUppercase'd accents
% issue: #63
\documentclass{report}
\usepackage[spanish]{babel}
\pagestyle{headings}
\begin{document}
\tableofcontents
\chapter{Introducción}
Texto con acción y canción.
\section{Antecedentes}
La teoría de la información estudia la cuantificación, el almacenamiento y la comunicación de la información digital. Fue establecida por Claude Shannon en los años cuarenta. La teoría de la información estudia la cuantificación, el almacenamiento y la comunicación de la información digital. Fue establecida por Claude Shannon en los años cuarenta. La teoría de la información estudia la cuantificación, el almacenamiento y la comunicación de la información digital.
\chapter{Conclusión}
Fin.
\begin{thebibliography}{9}
\bibitem{a} Un ítem.
\end{thebibliography}
\end{document}
